\documentclass[12pt]{article}
\usepackage{fancyhdr}
\pagestyle{fancy}
\fancyhf{}
\fancyfoot[C]{\thepage}
\fancyfoot[L]{\scriptsize CC BY-NC-SA 4.0 \textbullet\ Leo Liang}
\renewcommand{\headrulewidth}{0pt}
\usepackage[margin=1in]{geometry}
\usepackage{amsmath, amssymb}
\usepackage{array, booktabs}
\usepackage{pgffor}


\parindent=0pt
\parskip=1em
\newcolumntype{C}[1]{>{\centering\arraybackslash}p{#1}}

% writing lines
\newcommand{\wl}{\vspace{5mm}\noindent\rule{\linewidth}{0.4pt}\par\vspace{1mm}}
\newcommand{\wls}[1]{\foreach \i in {1,...,#1}{\wl}}

% ─────────────────────────────────────────────────────────────────────────────
\begin{document}
\pagestyle{fancy}

\begin{center}
  {\Large\bfseries Homework I}\\[4pt]
  Math Terminology $\cdot$ Speaking Math\\[6pt]
  Name:\;\underline{\hspace{6.5cm}}\quad
\end{center}
\hrule\bigskip

% =============================================================================
\section*{Instructions}
Complete the homework using \textbf{PEN}! If you have any questions or thoughts, write them down using a different color pen next to the problem/concept. Some homework questions are written intentionally vague, you should try to find them. Some information written may also contradict what you've already learned, you should try to find them. 

\section*{Part 1 --- Terminology Practice \hfill [100 pts]}

\subsection*{Foundation (45 pts)}

\textbf{1.}~[20 pts]\quad
Fill in the table for each term in $4x^3 - 3x^2y + 7x - 5$.

\medskip
\begin{center}
\renewcommand{\arraystretch}{2.2}
\begin{tabular}{|C{1.4cm}|C{2.3cm}|C{2.3cm}|C{2.5cm}|C{2.2cm}|}
\hline
\textbf{Term} & \textbf{Coefficient} & \textbf{Variable(s)}
              & \textbf{Exponent(s)} \\
\hline
$4x^3$   & & &  \\ \hline
$-3x^2y$ & & &  \\ \hline
$7x$     & & &  \\ \hline
$-5$     & & &  \\ \hline
\end{tabular}
\renewcommand{\arraystretch}{1}
\end{center}

\bigskip
\textbf{2.}~[15 pts]\quad
Underline each term. Write the number of terms in the blank.

\medskip
(a)\quad $3a^2b - a + 4b^3 - 1$
\hfill Number of terms:\;\underline{\hspace{2cm}}

\medskip
(b)\quad $x^2 + x + 1$
\hfill Number of terms:\;\underline{\hspace{2cm}}

\bigskip
\textbf{3.}~[10 pts]\quad
True or False. If false, write the corrected version on the line.

\medskip
(a)\quad T\,/\,F\quad $x$ has a coefficient of 1. \hrulefill

(b)\quad T\,/\,F\quad The expression $-15$ has no terms.
\hrulefill

(c)\quad T\,/\,F\quad A constant is not a term.
\hrulefill

(d)\quad T\,/\,F\quad $12x^3$ and $12x^2$ are like terms.
\hrulefill

% ─────────────────────────────────────────────────────────────────────────────
\subsection*{Application (35 pts)}

\textbf{4.}~[10 pts]\quad
Read the expression below out loud as one complete English sentence.
Write what you would say.
\[
  (x+2)(x-1)^2 + 12
\]

\wls{3}

\bigskip
\textbf{5.}~[15 pts]\quad
Write an algebraic expression for each description.

\medskip
(a)\quad ``Five times a number, squared''

\wl

(b)\quad ``Three terms: negative four times $y^3$, twice $y$, and one''

\wl

(c)\quad ``A number times itself, minus the same number''

\wl

\clearpage 
\bigskip
\textbf{6.}~[10 pts]\quad
A classmate says: \textit{``In $x^2 + x + 1$, the $x$-term has no
coefficient --- there's no number written in front of it.''}
Are they right? Explain.

\wls{2}


% =============================================================================
\section*{Part 2 --- Speaking: How to Say Algebra Out Loud}

\textit{Read this section carefully.
You will write expressions in English at the end.}

\bigskip
You can write algebra faster than you can say it. That's fine --- but if you
can't say an expression out loud, you don't fully own it yet. Reading your work
aloud is one of the best ways to catch mistakes, and every math class from here
on will expect you to communicate verbally. This section teaches you the
conventions that mathematicians actually use.

\bigskip
\noindent\textbf{Rule 1 --- Saying Powers}

Single-variable powers: $x^2$ = \textit{``x squared,''} $x^3$ = \textit{``x cubed,''}
$x^4$ = \textit{``x to the fourth.''} When there is a coefficient, say it first:
$3x^2$ = \textit{``three x squared.''} Coefficient, then variable, then power ---
same order as left to right.

\bigskip
\noindent\textbf{Rule 2 --- The Word ``Quantity''}

Whenever you see a \textbf{group} --- something inside parentheses --- you say
\textbf{``the quantity''} before it.
\[
  (x+2) \;=\; \textit{``the quantity } x \textit{ plus two''}
\]
``The quantity'' tells your listener: what I'm about to say is grouped together.
Think of it as speaking the open parenthesis in words.
\[
  (x+2)(x+1) \;=\;
  \textit{``the quantity } x \textit{ plus two, times the quantity }
  x \textit{ plus one''}
\]
Pay attention to where the comma is! This is crucial in both speaking and writing.

\bigskip
\noindent\textbf{Rule 3 --- Powers on Groups}

When an exponent sits on a whole group, say the group first, then the power at
the end. The exponent attaches to the group immediately in front of it.
\begin{align*}
  (x+1)^2 &\;=\; \textit{``the quantity } x \textit{ plus one, squared''} \\[4pt]
  (x+2)(x+1)^2 &\;=\;
    \textit{``the quantity } x \textit{ plus two, times the quantity }
    x \textit{ plus one, squared''} 
\end{align*}
\textbf{Note:} ``squared'' attaches to $(x+1)$ only. If the entire expression
were squared you would say ``\ldots\,all squared.'' Like this:
\begin{align*}
    \bigl[(x+2)(x+1)\bigr]^2 &\;=\;
    \textit{``the quantity } x \textit{ plus two, times the quantity }
    x \textit{ plus one, all squared''}
\end{align*}

\bigskip
\noindent\textbf{Rule 4 --- Fractions}

Fractions are read as \textit{``numerator over denominator.''} If either part is
a group, it gets ``the quantity.''
\begin{align*}
  \frac{x}{x+1} &\;=\;
    \textit{``}x\textit{ over the quantity }x\textit{ plus one''} \\[6pt]
  \frac{2x+1}{x^2-1} &\;=\;
    \textit{``the quantity two }x\textit{ plus one, over the quantity }
    x\textit{ squared minus one''}
\end{align*}

\bigskip
\noindent\textbf{Putting It Together}

Let's read $(x+2)(x+1)^2+12$ step by step:

\medskip
\begin{center}
\begin{tabular}{r @{\quad$\to$\quad} l}
\toprule
\textbf{Piece} & \textbf{What you say} \\
\midrule
$(x+2)$   & \textit{``the quantity $x$ plus two''} \\
$\times$  & \textit{``times''} \\
$(x+1)^2$ & \textit{``the quantity $x$ plus one, squared''} \\
$+\;12$   & \textit{``plus twelve''} \\
\bottomrule
\end{tabular}
\end{center}

\medskip
\textbf{Full reading:}
\textit{``The quantity $x$ plus two, times the quantity $x$ plus one,
squared, plus twelve.''}

\medskip
One more:\quad
$\dfrac{x(x+3)}{x^2+1}$ is read as 
\textit{``$x$ times the quantity $x$ plus three, over the quantity
$x$ squared plus one.''}

\bigskip
\subsection*{Now You Try}

Write a complete English sentence for each expression.

\medskip
\textbf{1.}\quad $(x+1)^2$

\wls{2}

\textbf{2.}\quad $\dfrac{x}{x+1}$

\wls{2}

\textbf{3.}\quad $(2x+1)(x-3)$

\wls{2}

\textbf{4.}\quad $\dfrac{x+2x}{x^2}$

\wls{2}

\textbf{5.}\quad $\dfrac{x^2+1}{(x-1)^2}$

\wls{2}

\bigskip
\textbf{6. Reverse.}

(a)\quad ``The quantity $2x$ minus one, squared.''

\wl

(b)\quad ``The quantity $x$ plus three, over the quantity $x$ minus three.''

\wl

(c)\quad ``Negative the quantity $x$ squared plus one.''

\wl

(c)\quad ``Negative $x$ squared plus one.''

\wl

(d)\quad ``The quantity $x$ plus one, times the quantity $x$ minus one,
all squared, over $x$.''

\wl

(e)\quad ``The quantity $x$ plus $h$, squared, minus $x$ squared, all over $h$.''

\wl

\section*{Diagnostic Problems}
Factor the following expressions. You may want to use a piece of scratch paper.

\textbf{(a)}\quad $x^2 - 16$

\wl

\clearpage
\textbf{(b)}\quad $2x^2 + 5x - 3$

\wl

\textbf{(c)}\quad $x^4 - 81$

\wl

\textbf{(d)}\quad $(x^2 + 3x)^2 - 2(x^2 + 3x) - 8$

\wl

\textbf{(e)}\quad $x^4 + x^2 + 1$

\wl

\textbf{(f)}\quad $x^4 + 4$

\wl

%============================

\section*{Part 3 --- The Language of Math}

\textit{Read carefully. These words appear in every math class from here on.}

\bigskip
Math uses ordinary English words in precise ways. When a problem says
\textit{evaluate}, it means something different from \textit{simplify}.
When you \textit{define} a variable, you are making a promise about what
that symbol means for the rest of your work. Here are the words you will
see most often:

\bigskip
\noindent\textbf{Let} --- introduces a variable.
``Let $x$ be the width of the room'' means: for everything that follows,
$x$ stands for that width. You are naming an unknown.

\bigskip
\noindent\textbf{Define} --- give an exact, unambiguous meaning to a symbol or term.
Defining $x = 5$ means $x$ cannot also be 7. A definition is a commitment.
If a problem asks you to \textit{define a variable}, respond with a
\textbf{let} statement: ``Let $x$ be the width of the room.''

\bigskip
\noindent\textbf{Given} --- a fact that has been established or handed to you as true.
You did not choose it; the problem did.
\textit{``Given $x=4$, find $2x$.''} means: this equation is true --- now work from it.

\bigskip
\noindent\textbf{Suppose/Assume} --- temporarily accept something as true
in order to explore what follows.
``Suppose the width is 10'' does not mean it is 10 --- it means
\textit{what would happen if it were?} 

\noindent\textbf{Suppose/Assume \ldots\ Then} --- temporarily accept something
as true in order to explore what follows. \textit{Then} delivers the consequence.
``Suppose the width is 10 blocks --- then the total area is $5(10+3) = 65$.''
The \textit{suppose} sets up the condition; the \textit{then} is what you conclude from it.
This is the skeleton of almost every mathematical argument you will ever write.

\begin{center}
    \begin{tabular}{@{} p{5cm} p{8cm} @{}}
    \toprule
    \textbf{Sentence} & \textbf{What it signals} \\
    \midrule
    \textit{Given that it is raining, bring an umbrella.}
      & It is actually raining. This is a known fact. Act on it. \\[6pt]
    \textit{Suppose it rains tomorrow, then bring an umbrella.}
      & We do not know if it will rain. We are planning for the possibility. \\[6pt]
    \textit{Assume the store is open, then we can buy wool tonight.}
      & We are temporarily treating this as true to figure out what to do next.
        It might be wrong. \\
    \bottomrule
    \end{tabular}
\end{center}

\bigskip
\noindent\textbf{Evaluate} --- substitute a value and compute a number.
``Evaluate $3x + 1$ when $x = 4$'' means replace $x$ with 4 and get 13.
The answer is always an expression.

\bigskip
\noindent\textbf{Simplify} --- rewrite an expression in a cleaner equivalent form.
The answer is still an expression, not necessarily a number.
$3x + 2x$ simplifies to $5x$; you have not changed the value, only the form.

\bigskip
\noindent\textbf{Determine} --- find with certainty, and show why.
Stronger than ``guess'' or ``find'' --- it implies your answer must be justified. An example questions could be, ``Determine the value of $x$ such that $2x + 3 = 11$.''

\bigskip
\noindent\textbf{Such that} --- a condition or constraint.
``Find $x$ such that $2x = 8$'' means find the $x$ that makes that statement true. Notice we could've also asked ``Determine $x$ such that $2x = 8$.''

\bigskip
\noindent\textbf{Hence / Therefore / Thus} --- what comes next follows
logically from what just came before.
These are not filler words. Writing ``therefore'' is a claim that
the next line is forced by the previous one. An example would be ``$2x = 8$, therefore $x = 4$.'' 

\bigskip
\noindent\textbf{Conjecture} --- an educated guess you believe is true but have not yet proven.
Mathematicians make conjectures all the time. Some stay unproven for centuries. You could be the next person to prove an important conjecture :)

\bigskip
\hrule
\bigskip
\noindent\textit{Quick check --- answer in one sentence each:}

\medskip
\textbf{1.}\quad What is the difference between \textit{evaluate} and \textit{simplify}?

\wls{2}

\textbf{2.}\quad A classmate writes ``let $x = 5$'' at the top of their work,
then uses $x = 7$ two lines later. What rule did they break?

\wls{2}

\textbf{3.}\quad Write a sentence using the word \textit{such that} about anything ---
it does not have to be about math.

\wls{2}

\textbf{4.}\quad Rewrite this using the word \textit{therefore}:
``It is raining. I should bring an umbrella.''

\wls{2}

\newpage
\textbf{5.}\quad Write a \textit{suppose \ldots\ then} sentence about anything.

\wls{2}

\textbf{6.}\quad Write a conjecture about even numbers.
You do not need to prove it.

\wls{2}

\textbf{7.}\quad Your friend says: ``I assumed the store was open,
so I walked over but found out it was actually closed.''
Did your friend use \textit{assumed} correctly? Explain.

\wls{2}


\textbf{8.}\quad Your friend says: ``I defined $n$ as an even number,
but then I realized it was odd, so I changed it in the middle of the problem.''
What rule did they break?

\wls{2}

% =============================================================================
\newpage
\section*{Part 4 --- Real Life Application}

You and a friend are building two bedrooms in Minecraft and need to
figure out how much wool to collect to carpet both rooms.
Each floor block needs exactly one carpet block, and each carpet block
takes two wool to craft.

\medskip
Both rooms are 10 blocks long and share a wall.
Bedroom~2 will be 8 blocks wide.
You haven't decided how wide Bedroom~1 will be yet.

\medskip
\textbf{1.}\quad How many carpet blocks do you need for Bedroom~2? \hrulefill

\textbf{2.}\quad How much wool do you need for Bedroom~2? \hrulefill

\textbf{3.}\quad How many carpet blocks do you need for Bedroom~1? \hrulefill

\textbf{4.}\quad How much wool do you need for Bedroom~1? \hrulefill

\textbf{5.}\quad How much wool do you need in total for both rooms? \hrulefill


\medskip
$\bigstar$\;\textit{Make sure to take notes in a different color pen of where your brain might be setting warning lights off!}



% =============================================================================
\newpage

\newpage
\begin{center}
  {\large\bfseries ANSWER KEY --- DO NOT DISTRIBUTE}\\[4pt]
  Homework 1 --- Math Terminology / Speaking Math
\end{center}
\hrule\bigskip

% ─────────────────────────────────────────────────────────────────────────────
\section*{Part 1 --- Terminology Practice}

\textbf{Q1.}
\begin{center}
\renewcommand{\arraystretch}{1.6}
\begin{tabular}{|C{1.4cm}|C{2.3cm}|C{2.3cm}|C{2.5cm}|C{2.2cm}|}
\hline
\textbf{Term} & \textbf{Coefficient} & \textbf{Variable(s)}
              & \textbf{Exponent(s)} & \textbf{Type} \\
\hline
$4x^3$   & $4$  & $x$    & $3$                    & variable \\
\hline
$-3x^2y$ & $-3$ & $x,y$  & $2$ on $x$; $1$ on $y$ & variable \\
\hline
$7x$     & $7$  & $x$    & $1$                    & variable \\
\hline
$-5$     & $-5$ & none   & none (or $x^0$)        & constant \\
\hline
\end{tabular}
\renewcommand{\arraystretch}{1}
\end{center}

\medskip
\textbf{Q2.}\quad
(a)~4 terms: $\underline{3a^2b}$, $\underline{-a}$, $\underline{4b^3}$, $\underline{-1}$\quad
(b)~3 terms: $\underline{x^2}$, $\underline{x}$, $\underline{1}$

\textit{Note: $-a$ has coefficient $-1$. The minus sign belongs to the term, not the separator.}

\medskip
\textbf{Q3.}\quad
(a)~\textbf{TRUE} --- $x = 1\cdot x$, coefficient is 1 by convention.\\
(b)~\textbf{FALSE} --- $-15$ is a constant term. It has one term.\\
(c)~\textbf{FALSE} --- a constant is a term (it has a coefficient and no variable).\\
(d)~\textbf{FALSE} --- like terms require identical variable parts; $x^3 \neq x^2$.

\medskip
\textbf{Q4.}\quad
\textit{``The quantity $x$ plus two, times the quantity $x$ minus one, squared, plus twelve.''}

Key points: ``the quantity'' before each grouped expression; ``squared'' attaches to $(x-1)$ only, not the whole product.

\medskip
\textbf{Q5.}\quad
(a)~$5x^2$ \textbf{or} $(5x)^2$ --- genuinely ambiguous. Both are defensible.
Good Socratic target: ask them which they chose and why.\\
(b)~$-4y^3 + 2y + 1$\\
(c)~$x^2 - x$

\medskip
\textbf{Q6.}\quad
No. The coefficient of $x$ is $1$ because $x = 1 \cdot x$.
A coefficient of 1 is invisible by convention, not absent.
Students must explain the \textit{why}, not just assert ``it's 1.''

% ─────────────────────────────────────────────────────────────────────────────
\section*{Part 2 --- Speaking Math}

\textbf{Now You Try:}

\textbf{1.}\quad \textit{``The quantity $x$ plus one, squared.''}

\textbf{2.}\quad \textit{``$x$ over the quantity $x$ plus one.''}

\textbf{3.}\quad \textit{``The quantity two $x$ plus one, times the quantity $x$ minus three.''}

\textbf{4.}\quad \textit{``The quantity $x$ plus two $x$, over $x$ squared.''}
Accept \textit{``three $x$ over $x$ squared''} if they simplified $x + 2x$ first.

\textbf{5.}\quad \textit{``The quantity $x$ squared plus one, over the quantity $x$ minus one, squared.''}

\medskip
\textbf{Reverse (Q6):}

(a)\quad $(2x-1)^2$

(b)\quad $\dfrac{x+3}{x-3}$

(c)\quad $-(x^2+1)$

(c)\quad $-x^2 + 1$\quad
\textit{Key distinction from the line above: no parentheses means the negative only applies to $x^2$.}

(d)\quad $\dfrac{\bigl[(x+1)(x-1)\bigr]^2}{x}$

(e)\quad $\dfrac{(x+h)^2 - x^2}{h}$\quad
\textit{This is the difference quotient --- the foundational expression of calculus.
If either student recognizes it or asks about it, tell them what it becomes.}

% ─────────────────────────────────────────────────────────────────────────────
\section*{Diagnostic --- Factoring}

\textbf{(a)}\quad $x^2 - 16 = (x+4)(x-4)$\quad [difference of squares]

\textbf{(b)}\quad $2x^2 + 5x - 3 = (2x - 1)(x + 3)$

\textit{Method: find two numbers that multiply to $2\times(-3)=-6$ and add to $5$: those are $6$ and $-1$.
Then split: $2x^2 + 6x - x - 3 = 2x(x+3) - 1(x+3) = (2x-1)(x+3)$.}

\textbf{(c)}\quad $x^4 - 81 = (x^2+9)(x^2-9) = (x^2+9)(x+3)(x-3)$

\textit{Apply difference of squares twice. Note $x^2+9$ does not factor further over the reals.}

\textbf{(d)}\quad $(x^2+3x)^2 - 2(x^2+3x) - 8$

Let $u = x^2 + 3x$:
\[
u^2 - 2u - 8 = (u-4)(u+2) = (x^2+3x-4)(x^2+3x+2)
\]
Factor each quadratic:
\[
= (x+4)(x-1)(x+2)(x+1)
\]

\textit{Students who get here have strong Algebra 1 instincts. The substitution step is the key insight.}

\textbf{(e)}\quad $x^4 + x^2 + 1$

Add and subtract $x^2$:
\[
x^4 + x^2 + 1 = x^4 + 2x^2 + 1 - x^2 = (x^2+1)^2 - x^2 = (x^2+x+1)(x^2-x+1)
\]

\textbf{(f)}\quad $x^4 + 4$\quad [Sophie Germain Identity]

Add and subtract $4x^2$:
\[
x^4 + 4 = x^4 + 4x^2 + 4 - 4x^2 = (x^2+2)^2 - (2x)^2 = (x^2+2x+2)(x^2-2x+2)
\]

\textit{Neither quadratic factors further over the reals.
If a student cracks either (e) or (f), they are thinking at a competition level.}

% ─────────────────────────────────────────────────────────────────────────────
\section*{Part 3 --- Language of Math (Quick Check)}

\textbf{1.}\quad \textit{Evaluate} means substitute a specific value and produce a \textbf{number}.
\textit{Simplify} means rewrite in a cleaner equivalent form --- the result is still an \textbf{expression}.

\textit{Note: there is a typo in the reading --- it says ``evaluate always gives an expression.''
It should say ``number.'' See if they catch it. (This is one of the intentional contradictions flagged in the instructions.)}

\textbf{2.}\quad They broke the rule that \textbf{let} (and \textbf{define}) is a commitment.
Once you write ``let $x = 5$,'' $x$ is 5 for the entire problem. You cannot change it.

\textbf{3.}\quad Open-ended. Any grammatically correct sentence using ``such that'' as a constraint.
\textit{Example: ``Find a number such that doubling it gives 10.''} 
Watch for students who use it as decoration rather than as a condition.

\textbf{4.}\quad \textit{``It is raining, therefore I should bring an umbrella.''}

\textbf{5.}\quad Open-ended. Look for correct structure: a condition followed by a consequence.
\textit{Example: ``Suppose school is cancelled, then I can sleep in.''}

\textbf{6.}\quad Open-ended. Any claim about even numbers not yet proven.
\textit{Strong examples: ``The sum of any two even numbers is even.''
``Every even number greater than 2 is the sum of two primes.'' (That one is Goldbach's Conjecture.)}

\textbf{7.}\quad \textbf{Yes}, they used it correctly.
\textit{Assume} means accepting something as true without certainty in order to act.
It turned out to be wrong --- and that is fine. \textit{Assume} does not promise correctness, only temporary acceptance.

\textbf{8.}\quad They broke the rule that \textbf{define} is a commitment.
Once $n$ is defined as even, it cannot become odd mid-problem.
If they needed $n$ to be odd, they should have defined it that way from the start or used a different variable.

% ─────────────────────────────────────────────────────────────────────────────
\section*{Part 3 --- Real Life Application}

\textbf{1.}\quad $10 \times 8 = \mathbf{80}$ carpet blocks

\textbf{2.}\quad $80 \times 2 = \mathbf{160}$ wool

\textbf{3.}\quad Cannot be a number --- the width is unknown.
Accept any representation: $10w$, $10 \times ?$, ``$10$ times the width.''
The key move: did they \textbf{name} the unknown? Award full credit only if they introduce a symbol or label.

\textbf{4.}\quad $2 \times 10w = \mathbf{20w}$ wool (using whatever variable they defined)

\textbf{5.}\quad $20w + 160$ wool total

\textit{Diagnostic note: a student who writes $20w + 160$ without prompting
has independently constructed a linear function. That is the bridge to Lesson 2.}

\end{document}
